\section*{الفصل \ref{ch:b3:developmental-genetics} --- وراثة النماء وتطور الشكل}

\begin{solution}{exo:b3:developmental-genetics:1}
الأثر الأمومي: \emph{bicoid}، فأجنّة الأمهات الطافرات تفتقر
إلى الرأس والصدر. والفجوية: \emph{Krüppel}، فكتلة متصلة من
القطع مفقودة. والزوجية القاعدة: \emph{fushi tarazu} (أو
\emph{even-skipped})، فقطعة بعد أخرى مفقودة. وقطبية القطعة:
\emph{hedgehog} (أو \emph{engrailed})، فجزء من كل قطعة
مستبدَل بصورة مرآتية للباقي.
\end{solution}

\begin{solution}{exo:b3:developmental-genetics:2}
مورّثة يودع منتجها في البيضة الأمُّ، فيتبع النمط الظاهري
للجنين نمطها الوراثي. وأجنّتها تفتقر إلى الرأس والصدر وتحمل
بنى خلفية عند الطرفين؛ وهي تموت. والأليل الأبوي البري في
اللاقحة، غير أن رسول \emph{bicoid} تصنعه خلايا الأم المرضعة
أثناء تكوّن البيوض ويُثبَّت عند القطب قبل الإلقاح؛ ونسخة
اللاقحة نفسها كانت ستُستنسخ بعد ساعات، بعد أن يكون وقت
التدرج قد مضى.
\end{solution}

\begin{solution}{exo:b3:developmental-genetics:3}
\omterm{def:b3:developmental-genetics:hox}{التوافق الخطي}: ترتيب \omterm{def:b3:developmental-genetics:hox}{مورّثات Hox} على امتداد الصبغي يطابق
ترتيب نطاقات تعبيرها على امتداد الجسم. والغلبة الخلفية: حيث
يُعبَّر عن عدة منها، تحدد الأكثر خلفيةً الهويةَ. والفأر
الذي يفتقر إلى \emph{Hoxc8} يُظهر تحولًا أماميًا --- إذ تتخذ
فقرته القطنية الأولى هوية فقرة صدرية وتحمل ضلعًا. و\emph{Antennapedia}
في الرأس يحوّل قرون الاستشعار إلى أرجل، وهي هوية القطعة
الصدرية الثانية.
\end{solution}

\begin{solution}{exo:b3:developmental-genetics:4}
شفة ظهرية مطعَّمة استحثت محورًا ثانيًا كاملًا مصنوعًا في
معظمه من خلايا المضيف: فقد أمر الطعم جيرانه --- أي \omterm{def:b3:developmental-genetics:organiser}{منظِّم}.
وجزيئاته مضادات مفرَزة لبروتين BMP (الكوردين، والنوغين،
والفوليستاتين). والأديم الظاهر يصير عصبيًا حين تغيب إشارة
BMP وبشرةً حين تحضر، فالعصبي هو ما يفعله الأديم الظاهر حين
لا يُخبَر بشيء: \omterm{def:b3:developmental-genetics:organiser}{فالمنظِّم} يعمل بإسكات أمر.
\end{solution}

\begin{solution}{exo:b3:developmental-genetics:5}
$k = \ln 2/(35\times 60) = \qty{3.3e-4}{s^{-1}}$؛ و$\lambda = \sqrt{4/
3.3\times 10^{-4}} = \qty{110}{\micro m}$؛ والنسبة $\mathrm{e}^{250/110}
= \mathrm{e}^{2.27} = 9.7$.
\end{solution}

\begin{solution}{exo:b3:developmental-genetics:6}
أربع نسخ تضاعف $c_{0}$: فينتقل الحد إلى الوراء بمقدار $\lambda\ln 2 =
\qty{69}{\micro m}$ إلى \qty{309}{\micro m}؛ ونسخة واحدة تنصّفه:
فينتقل إلى الأمام بمقدار \qty{69}{\micro m} إلى \qty{171}{\micro m}.
وكل إزاحة \qty{14}{\%} من طول البيضة --- أكبر مما يُرصد، وهو
نفسه شاهد على آليات تخمد اعتماد الجرعة.
\end{solution}

\begin{solution}{exo:b3:developmental-genetics:7}
$\delta x = \lambda\,\delta c/c$: فلأجل \qty{8}{\micro m} يكون $\delta c/c =
8/100 = \qty{8}{\%}$. والعدّة البواسونية $N$ خطؤها
النسبي $1/\sqrt{N}$، فيجب عدّ $N = 1/0.08^{2} \approx 160$ جزيئًا
--- أي بضع في المئة من الجزيئات الحاضرة في نواة، أو عدّة
واحدة مكاملة على الزمن.
\end{solution}

\begin{solution}{exo:b3:developmental-genetics:8}
$k = \ln 2/35 = \qty{0.0198}{min^{-1}}$: فتعطي $1 - \mathrm{e}^{-kt}$
قيمة $0.70$ عند \qty{60}{min}، و$0.83$ عند \qty{90}{min}، و$0.91$ عند
\qty{120}{min}، و$0.95$ عند \qty{150}{min}. فالتدرج لا يزال يرتفع
بضع في المئة وهو يُقرأ، وهو ما كان سيزيح حدًّا بمقدار
$\lambda\ln(1/0.9) \approx \qty{10}{\micro m}$، أي نحو نواة، على
مدى نافذة القراءة: فعلى الجنين إما أن يقرأ في زمن ثابت وإما
أن يستعمل قراءة غير حساسة للمستوى الكلي.
\end{solution}

\begin{solution}{exo:b3:developmental-genetics:9}
منطقة نشاط مستقطِب ثانية تعطي مضاعفةً مرآتية، 4-3-2-2-3-4:
إذ ترى الخلايا الأمامية الآن Shh مرتفعًا فتصير أصابع خلفية.
وحبيبة تطلق قليلًا من Shh تعطي مضاعفةً جزئية --- إصبعًا 2
أماميًا إضافيًا، مثل 2-2-3-4 --- لأن الخلايا الأمامية لا ترى
إلا التركيز المنخفض الذي يحدد الإصبع 2.
\end{solution}

\begin{solution}{exo:b3:developmental-genetics:10}
الحل الخاص $s/k$؛ والحلان المتجانسان $\cosh(x/\lambda)$
و$\sinh(x/\lambda)$؛ وانعدام الدفق عند $0$ يقتل $\sinh$؛ و$c(L) = 0$
يحدد السعة: $c(x) = (s/k)\bigl[1 - \cosh(x/\lambda)/\cosh(L/\lambda)\bigr]$.
والمقطع هضبة قرب $x = 0$ تهبط إلى الصفر في طبقة حدية
عرضها $\lambda$ عند البالوعة --- لا أسّيًا من نقطة، بل حقلًا
مستويًا بحافة؛ وأعلى تركيز أبعد ما يكون عن البالوعة. وعند
$L \ll \lambda$ يكون $\cosh u \approx 1 +
u^{2}/2$ و$c \approx s(L^{2} - x^{2})/(2D)$: أي قطع مكافئ مستقل
عن $k$، إذ تبلغ الجزيئات البالوعة قبل أن تقدر على الاضمحلال.
\end{solution}

\begin{solution}{exo:b3:developmental-genetics:11}
$x_{1} = \lambda\ln(1/0.3) = 1.20\lambda$ و$x_{2} = \lambda\ln(1/0.08)
= 2.53\lambda$: فعرض المنطقة الأولى $1.20\lambda$، والثانية
$1.32\lambda$، وكلاهما مستقل عن $c_{0}$، الذي يزيحهما معًا.
والمنطقة الثالثة، $L - 2.53\lambda$، تمتص كل التباين في طول
البيضة: ففي بيضة \qty{450}{\micro m} عند $\lambda = \qty{100}{\micro m}$
يكون البطن \qty{197}{\micro m}، وفي بيضة \qty{550}{\micro m}
\qty{297}{\micro m} --- فالنمط غير مدرَّج. والآليات: تدرج ثانٍ
من القطب الخلفي (Nanos، أو Caudal، أو الجهاز الطرفي) يُقرأ مع
Bicoid، فتتحدد المواضع بالنسبة؛ أو جزيء ``موسِّع'' يُصنع حيث
يكون المورفوجين منخفضًا فيطيل $\lambda$ حتى يملأ التدرج
البيضة.
\end{solution}

\begin{solution}{exo:b3:developmental-genetics:12}
التغيّر الترميزي يبدّل البروتين حيثما عمل: فالمورّثة
\emph{Pitx1} تبني \omterm{def:b3:endocrinology:axes}{النخامى} والفك أيضًا، و\emph{\omterm{def:b3:developmental-genetics:organiser}{sonic hedgehog}}
ينمّط الأنبوب العصبي والمعي والأسنان أيضًا، فالبروتين المكسور
قاتل أو معطِّل بتعدد الأثر. وأما المعزِّز فيسوق التعبير في
نسيج واحد في زمن واحد؛ وحذفه يزيل البنية التي كان ذلك النسيج
سيبنيها ويترك كل استعمال آخر للبروتين سليمًا. والمعزِّزات
نميطية وكثيرًا ما تكون فائضة جزئيًا، فتكون الخطوات الوسيطة
قابلة للحياة، وقد سُلك الطريق نفسه --- فقدُ مفتاح --- استقلالًا
في بحيرات كثيرة من بحيرات شوكيات الظهر.
\end{solution}

\begin{solution}{pb:b3:developmental-genetics:1}
\textbf{1.} $k = \ln 2/(35\times 60) = \qty{3.3e-4}{s^{-1}}$؛ و$1/k =
\qty{3030}{s} = \qty{50}{min}$.
\textbf{2.} $\lambda = \sqrt{4/3.3\times 10^{-4}} = \qty{110}{\micro m}$؛
والبيضة $4.5\lambda$.
\textbf{3.} $\mathrm{e}^{4.5} \approx 90$ من قطب إلى قطب؛
و$\mathrm{e}^{8.5/110} = 1.08$، أي تغيّر \qty{8}{\%} لكل نواة.
\textbf{4.} $J = 50\times\sqrt{4\times 3.3\times 10^{-4}} = 50\times
0.036 = \qty{1.8}{nM\,\micro m\,s^{-1}}$. والنانومولار الواحد $0.6$ جزيء
في \unit{\micro m^3}، فنحو $1.1$ جزيء في \unit{\micro m^2}
في الثانية.
\textbf{5.} حجم النواة $\tfrac{4}{3}\pi\times 27 = \qty{113}{\micro
m^3}$. وعند \qty{50}{nM}، $30$ جزيئًا في \unit{\micro m^3}:
أي \num{3400} جزيء. وعند منتصف البيضة $c = 50\,\mathrm{e}^{-2.27} =
\qty{5.2}{nM}$: أي نحو \num{350} جزيء.
\textbf{6.} $kt = 1.19$ عند \qty{60}{min}: $0.70$؛ و$kt = 2.38$ عند
\qty{120}{min}: $0.91$. فالتدرج في حدود \qty{10}{\%} من المستقر
حين يُقرأ --- ولا يزال يرتفع بمقدار كان سيحرّك حدًّا بمقدار
$\lambda\ln(1/0.91) \approx \qty{10}{\micro m}$، أي نواة.
\textbf{7.} $x_{1} = 110\ln(1/0.3) = \qty{132}{\micro m}$
(\qty{26}{\%})؛ و$x_{2} = 110\ln(1/0.08) = \qty{278}{\micro m}$
(\qty{56}{\%}).
\textbf{8.} أربع نسخ: $c_{0} = \qty{100}{nM}$؛ وينتقل الحدان
إلى الوراء بمقدار $\lambda\ln 2 = \qty{76}{\micro m}$، إلى $208$ و
\qty{354}{\micro m}. وست نسخ: $\lambda\ln 3 = \qty{121}{\micro m}$،
إلى $253$ و\qty{399}{\micro m}.
\textbf{9.} نسخة واحدة: إلى الأمام بمقدار \qty{76}{\micro m}، إلى $56$ و
\qty{202}{\micro m}. فمنطقة الرأس صغرت من $132$ إلى
\qty{56}{\micro m}؛ والصدر حفظ عرضه (\qty{146}{\micro m}) وانتقل
إلى الأمام؛ والبطن كبر من $222$ إلى \qty{298}{\micro m}.
\textbf{10.} $\delta x = 110\times 0.1 = \qty{11}{\micro m}$، أي $1.3$
تباعد نووي.
\textbf{11.} نواة واحدة: $\delta c/c = 8.5/110 = \qty{7.7}{\%}$. ومن
\qty{10}{\%}، $n = (10/7.7)^{2} \approx 1.7$: أي قراءتان مستقلتان
(أو متوسط نواتين متجاورتين) يكفيان.
\textbf{12.} $1/\sqrt{350} = \qty{5.3}{\%}$: فالعدّة الآنية الواحدة
أدق أصلًا من \qty{10}{\%} المفترضة، فالضجيج المرصود ليس ضجيج
عدّ وحده بل يأتي من حركيات الارتباط ومن آلة القراءة ---
والتوسيط على الزمن يحسّن عليه.
\textbf{13.} $278/450 = \qty{62}{\%}$ و$278/550 = \qty{51}{\%}$ من
طول البيضة: أي تفاوت \qty{11}{\%}، ست نُوى. فالنمط لا يتدرج
عند $\lambda$ ثابت.
\textbf{14.} إذا كان $D \propto L^{2}$ مع ثبات $k$، كان $\lambda \propto L$
و$x_{T}/L = (\lambda/L)\ln(c_{0}/T)$ نفسه في كل بيضة: أي تدريج
تام. غير أن $D$ خاصية للجزيء وللعصارة، لا لحجم البيضة، فذلك
غير معقول كما صيغ؛ والتدريج في الأجنّة الحقيقية جزئي ويأتي
من قراءات أخرى.
\textbf{15.} $c^{5}/(c^{5} + T^{5}) = 0.1$ عند $c/T = 9^{-1/5} = 0.64$ و
$0.9$ عند $c/T = 9^{1/5} = 1.55$: أي عامل $2.4$ في التركيز،
و$\Delta x = 110\ln 2.4 = \qty{97}{\micro m}$ --- أي إحدى عشرة نواة،
أوسع بكثير من الحد المرصود البالغ نواتين أو ثلاثًا؛ فقراءة
التدرج وحدها ليست حادة بما يكفي.
\textbf{16.} الحدّة تُسقط تركيزًا معينًا على تشغيل أو إطفاء
إسقاطًا أحسم، فتضيق منطقة الانتقال؛ غير أن خطأً قدره
\qty{10}{\%} في التركيز لا يزال يحرّك النقطة التي تُعبَر
عندها العتبة بمقدار $\lambda\,\delta c/c$: فالاستجابة قد تكون
درجة تامة وتقع الدرجة مع ذلك في الموضع الخطأ.
\textbf{17.} لو كانتا مشغَّلتين معًا لكانت كل واحدة تكبت
الأخرى، فتُدفع إحداهما إلى الإطفاء على الأقل؛ ومع الكبت
التعاوني تكون الحالتان المستقرتان الوحيدتان واحدةً مشغَّلة
وأخرى مطفأة، ويكون الحد بينهما مفتاحًا لا ميلًا.
\textbf{18.} $\mathrm{d}\ln\lambda = \tfrac{1}{2}(0.02 - 0.06) = -0.02$
لكل درجة: فيهبط $\lambda$ بمقدار \qty{2}{\%} لكل درجة، أي
\qty{10}{\%} عند \qty{5}{\celsius}، إلى \qty{99}{\micro m}؛ وينتقل
حدّ \emph{hunchback} من $278$ إلى \qty{250}{\micro m}، أي ثلاث
نُوى إلى الأمام. والذباب ينمو بين \qty{18}{} و\qty{29}{\celsius}
بنسب سوية، فلا بد أن شيئًا يعوّض --- منبعًا أو قراءةً تنزاح
مع الحرارة في الاتجاه الآخر.
\textbf{19.} $2^{8} = 256$ شفرة؛ وأربع عشرة مستعملة. ومع
الغلبة الخلفية لا تهم إلا المورّثة الأكثر خلفيةً المشغَّلة،
فتكون الهويات المتمايزة تسعًا على الأكثر (ثماني مورّثات، أو
لا شيء): فالشفرة ترتيب لا كلمة ثنائية.
\textbf{20.} تصير القطعة الصدرية الثالثة ثانيةً: فقد كانت
\emph{Ubx} تكبت برنامج الجناح الذي تشغّله قطعة صدرية
افتراضًا، وإزالة الكابت تطلق الحالة السلفية --- أي أربعة
أجنحة، كما عند أسلاف الذبابة ذوي الأجنحة الأربعة.
\textbf{21.} المنطقة الرقبية في الإوزة تمتد إلى الفقرة 22: أي
عنق طويل من فقرات كثيرة، في مقابل سبع في الفأر. والمورّثة
واحدة؛ وإنما انتقل الحد الأمامي لتعبيرها خلفيًا --- أي تغيّر
تنظيمي في أين تُشغَّل مورّثة Hox.
\textbf{22.} صارت كل حافة منبعًا؛ فيكون Shh أعلى ما يكون عند
الحافتين (الإصبع 4)، ويهبط نحو الوسط (3، ثم 2)، ولا يبلغ قط
المستوى المنخفض الذي يحدد الإصبع 1، فيحمل الوسط إصبعين 2
ظهرًا لظهر: 4-3-2-2-3-4.
\textbf{23.} \emph{Pax6} \omterm{def:b3:developmental-genetics:organiser}{منظِّم} سيد محفوظ يقدر على إطلاق أي
برنامج عين يحمله المضيف: فالذبابة بنت عين ذبابة، فمورّثة
الفأر لا تحمل معلومة عين فأر، بل الأمر ``ابنِ عينًا هنا'' لا
غير. وهو يدل على \omterm{def:b3:developmental-genetics:evodevo}{تماثل عميق} للشبكة المورّثية، لا على تماثل
العضوين: فالعين المركّبة وعين الكاميرا تطورتا منفصلتين من
سلف كانت له رقعة مستقبِلة ضوئية ذات \emph{Pax6}.
\textbf{24.} خمسة أصابع على امتداد \qty{2}{mm} طول موجة قدره
\qty{0.4}{mm}؛ وستة تحتاج إلى \qty{0.33}{mm}، أي تقصير بالسدس.
ولأن $\lambda \propto \sqrt{D_{\text{inh}}/k_{\text{inh}}}$، ارفع
اضمحلال المثبِّط \qty{44}{\%} أو اخفض انتشاره \qty{31}{\%}؛ وفي
صفيحة اليد يفعل ذلك خفضُ جرعة \omterm{def:b3:developmental-genetics:hox}{مورّثات Hox} الخلفية، والفئران
الطافرة لها ستة أصابع أو سبعة أو أكثر، رفيعة.
\textbf{25.} $\lambda \approx \qty{110}{\micro m}$؛ ومضاعفة المنبع
تزيح كل حد بمقدار \qty{76}{\micro m}؛ وضجيج \qty{10}{\%} في
التركيز خطأ موضعي قدره \qty{11}{\micro m}، أي نحو نواة واحدة.
\end{solution}
